\section{Related work}
\label{sec:related}

\paragraph{Reading intermediate layers.}
The logit lens~\citep{nostalgebraist2020interpreting} decodes an intermediate state through the model's output layer, and the tuned lens~\citep{belrose2023eliciting} learns a decoder for each layer. Both show predictions sharpening with depth. Future Lens~\citep{pal2023future} and Patchscopes~\citep{ghandeharioun2024patchscopes} read richer content from a single hidden state. All of these read the state through the output vocabulary. Our places are fitted from the states themselves, and reading them does not need the answer words' output directions.

\paragraph{How depth is used.}
Several studies find that layers play different roles. Early layers gather and detokenise, middle layers build abstract features, and late layers refine a prediction~\citep{lad2024the,csordas2025do,gupta2025how,cheng2024emergence,skean2025layer}. Correct and incorrect predictions follow similar confidence curves across layers~\citep{kim2025on}. In factual recall, information about the relation and then the subject moves to the last position, where attention heads extract the attribute~\citep{geva2023dissecting}. Feed-forward layers act as key--value memories~\citep{geva2020transformer}, and middle-layer feed-forward modules at the subject's last token are causally decisive~\citep{meng2022locating}. Our onset depth marks, for a whole population of questions, the depth at which each model's answer becomes available at the last position.

\paragraph{Representation geometry.}
Middle layers represent the same content in a shared space across languages and modalities~\citep{wendler2024do,wu2024the}. Concepts have structured geometry. Categories form simplices~\citep{park2024the}, some features are circular~\citep{engels2024not}, space and time are linearly decodable~\citep{gurnee2023language}, counts lie on curved manifolds that attention reshapes~\citep{gurnee2026when}, and truth has a linear direction~\citep{marks2023the}. In trained classifiers, representations of different classes separate by a roughly constant factor per layer~\citep{he2022a}. Language models move partway toward neural collapse, in which states gather at their class means, at their top layer, and further with scale~\citep{wu2024linguistic}. The residual stream contains stable regions, within which small changes leave the next-token prediction unchanged~\citep{janiak2024characterizing}. Yet, almost surely, no two prompts map to the same last-position state~\citep{nikolaou2025language}. Places fit both findings: questions can gather in the directions that carry the answer while staying distinct in others.

\paragraph{What models know but do not say.}
Internal states often encode more than the output shows~\citep{orgad2024llms,gekhman2025insideout,ruscio2026on}. On knowledge benchmarks such as MMLU and ARC, however, this gap is near zero~\citep{park2025bridging}, and internal states mostly reflect what the model recalls rather than what is true~\citep{cheang2026do}. False few-shot demonstrations can override a correct internal answer at a critical layer~\citep{halawi2023overthinking}, and users' suggestions shift answers~\citep{sharma2023towards}. In multiple choice, attention in a few middle layers selects the answer letter, and a sparse set of heads then raises its probability~\citep{wiegreffe2024answer}. The winning option is chosen by its content before it is bound to its letter~\citep{wong2026when}, and models carry a prior over letters~\citep{zheng2023large}.

\paragraph{Steering.}
Adding a direction to the residual stream changes a model's behaviour~\citep{panickssery2023steering,li2023inferencetime}. Our moves are of this kind. Each one, however, takes a single question's state from one measured place to another at one layer, and is compared with random moves of the same length.
